\chapter{Background and Related Work}
\label{ch:background}

\section{The Biology of Epileptic Seizures}

Before discussing how seizures are detected computationally, it is worth being explicit about what a seizure actually is physiologically, since the detection problem this dissertation addresses is a direct consequence of it. A seizure is a transient episode of abnormally excessive or synchronous neuronal electrical activity in the brain: a large population of cortical neurons that would ordinarily fire in a loosely correlated, low-amplitude fashion instead begins firing in a highly synchronised, high-amplitude, rhythmic pattern, typically triggered by a local disruption to the normal balance between excitatory and inhibitory signalling in a region of cortex \citep{who2024epilepsy}. This abnormal synchrony can remain localised (a focal seizure) or spread across both hemispheres (a generalised seizure), and clinicians conventionally divide a seizure's timeline into the \emph{preictal} period beforehand, the \emph{ictal} period during the seizure itself, and the \emph{postictal} period of recovery afterward -- terminology used throughout this dissertation, particularly in Chapter~\ref{ch:results}, where "ictal morphology" refers to the characteristic shape this synchronised electrical activity takes in a given patient's recording.

Electroencephalography (EEG) is able to detect this abnormal synchrony non-invasively because scalp electrodes measure the summed electrical field potential generated by large populations of cortical neurons firing together; a small number of neurons firing independently produces a signal too weak and too spatially cancelled-out to register at the scalp, but the highly synchronised firing of a seizure produces a comparatively large, coherent voltage deflection that does register, which is precisely why EEG -- rather than, say, single-neuron recording -- is the practical tool of choice for this problem at scale. This also explains why seizure detection is fundamentally a pattern-recognition problem over time-series waveform morphology rather than a simple amplitude-threshold problem: a seizure is identifiable by its characteristic rhythmic, evolving waveform shape, not merely by being "louder" than the surrounding background EEG, and a substantial part of what a detection model -- whether a hand-crafted-feature classifier or a learned convolutional network, as in this project -- must actually learn is this shape.

\realfig{eeg_seizure_vs_baseline_example_chb10.png}{0.85}{A side-by-side raw EEG waveform comparison from patient chb10: a short segment of genuine seizure activity next to a short segment of quiet interictal background, same channel and y-axis scale, illustrating the waveform morphology difference a detection model must learn.}{fig:eeg-waveform-example}

\realfig{eeg_seizure_vs_baseline_example_chb13.png}{0.85}{The same seizure-versus-baseline comparison for patient chb13 -- shown alongside chb10 above to illustrate that this waveform-morphology distinction, and the difficulty of learning it, varies noticeably from patient to patient.}{fig:eeg-waveform-example-chb13}

\section{Epileptic Seizure Detection from Scalp EEG}

Automated seizure detection from scalp electroencephalography (EEG) is a long-studied problem, and the CHB-MIT paediatric dataset \citep{goldberger2000} -- continuous, long-duration recordings from young patients with intractable epilepsy, collected at Boston Children's Hospital -- has become one of its standard public benchmarks. The dataset's realism is also its central difficulty: seizures are rare relative to interictal (non-seizure) background, recordings run for many hours per patient, and patients differ substantially from one another in electrode montage coverage, seizure semiology, and background EEG morphology.

Two evaluation conventions recur throughout the literature and are both used, for different purposes, in this dissertation. \emph{Window-level} evaluation scores every fixed-length segment of EEG independently, as an ordinary binary classification problem; it is convenient for training but does not directly correspond to what a clinician or patient cares about, since a single seizure event may span many consecutive windows, none of which need to be individually correct for the event as a whole to be usefully flagged. \emph{Event-level} evaluation instead asks whether each genuine seizure event was flagged at all (event sensitivity), and how often the system fires outside of any genuine event, expressed as a rate per hour of recording (FP/hr) rather than a raw count, so that patients with different total recording durations remain comparable. This dissertation treats event-level sensitivity and FP/hr as its primary reporting unit throughout, for the clinical reasons set out in Section~\ref{sec:eval-protocol}, while retaining window-level metrics as a diagnostic tool -- most importantly, window-level specificity as the basis of a \emph{collapse diagnostic} used throughout this dissertation. In brief, since it is used repeatedly from this point onward: a model that fires on almost every window it sees will register a high event sensitivity simply because it never stays silent through a genuine seizure, which is an artefact of near-permanent triggering rather than genuine discrimination; the collapse diagnostic is a gate, based on window-level specificity and predicted-block length, that catches this artefact before a sensitivity figure is reported as genuine. Its exact numerical criteria are given in full in Section~\ref{sec:eval-protocol}.

\section{Neuromorphic Computing and the BrainChip AKD1000}

Conventional deep learning inference performs a multiply-accumulate operation at every neuron, at every clock tick, regardless of whether the underlying input has changed -- a computational pattern well suited to GPU batching but poorly suited to a power budget measured in microwatts. Neuromorphic computing takes its name and its founding motivation from the alternative the biological brain demonstrates: sparse, event-driven computation, in which a neuron performs work only when it receives an incoming spike, and most neurons are silent at any given moment \citep{mead1990}. \citet{schuman2017} give a broad survey of how this principle has been implemented across a range of neuromorphic hardware platforms, of which BrainChip's AKD1000 is one commercial example.

The AKD1000 v1 implements this event-driven principle in silicon via an array of Neural Processing Units (NPUs) connected by an event-driven on-chip fabric, with synaptic weights stored in low-bitwidth on-chip SRAM rather than requiring external DRAM \citep{brainchip_akida_whitepaper, brainchip_akd1000_productbrief}. The architectural consequence that matters most for this project is that an NPU which receives no spike performs no dynamic computation and draws negligible power; on a sparsely-firing spiking network -- which, as Section~\ref{ch:results} shows, an EEG-derived signal typically produces, given how much of the recording is quiet interictal background -- the majority of the chip's NPUs are idle at any given moment, in direct contrast to a GPU, where every core is occupied for every tensor operation regardless of the input's information content.

This power efficiency is not free. The AKD1000 v1 imposes a specific, hard set of architectural constraints that shaped every design decision in this project's model architecture (Chapter~\ref{ch:methodology}), summarised in full in Table~\ref{tab:akd1000-constraints}. Every model architecture variant considered in this project was designed to satisfy these constraints from the outset, rather than being adapted to them after an incompatible design had already been trained; the full constraints reference document is maintained alongside the project's codebase at \url{https://github.com/samidR25/dissertation} for any constraint not covered by the summary below.

\begin{table}[H]
\centering
\caption{AKD1000 v1 architectural constraints relevant to this project's model design.}
\label{tab:akd1000-constraints}
\begin{tabular}{p{4cm}p{5cm}p{4.5cm}}
\toprule
Constraint & Value & Consequence \\
\midrule
Weight/activation bitwidth & Maximum 4 bits (w4a4) & Confirmed via \texttt{cnn2snn}'s own model-compatibility validation as this hardware revision's hard ceiling \\
\midrule
Dilated convolutions & Unsupported (\texttt{dilation\_rate} $\neq 1$) & Fails at \texttt{check\_model\_compatibility()}; excluded from every architecture variant, including longer-receptive-field designs considered early in the project \\
\midrule
Batch normalisation & Must be folded or removed prior to conversion & No direct spiking equivalent exists; addressed via the BN-placement fix discussed in Chapter~\ref{ch:methodology} \\
\midrule
Kernel sizes & Restricted to 1$\times$1, 3$\times$3, 5$\times$5, 7$\times$7 & Fixed silicon routing; larger or irregular kernels are not representable \\
\midrule
Activation function & ReLU only & The ANN-to-SNN conversion (Section~\ref{sec:ann-snn-conversion}) is mathematically defined for ReLU and does not extend to LeakyReLU, GELU, or similar variants \\
\midrule
Skip connections & Not supported in early design iterations & Residual paths complicate conversion; simplified out of every architecture variant used for reported results \\
\bottomrule
\end{tabular}
\end{table}

\section{ANN-to-SNN Conversion}
\label{sec:ann-snn-conversion}

The AKD1000 does not train spiking neural networks (SNNs) natively from scratch. Instead, the standard and the project's own workflow is to train a conventional artificial neural network (ANN) with ReLU activations using ordinary backpropagation, then convert it into an equivalent spiking form for deployment. \citet{maass1997} first articulated SNNs as a distinct, biologically grounded computational paradigm -- a "third generation" of neural network model, in which information is represented in the timing and rate of discrete spikes rather than continuous real-valued activations. \citet{gerstner2002} give the standard mathematical treatment of the underlying single-neuron model most commonly used in this conversion, the Leaky Integrate-and-Fire (LIF) neuron, which accumulates incoming current on a membrane potential that leaks toward rest over time and emits a spike when that potential crosses a threshold, resetting immediately afterward. Formally, the membrane potential $V(t)$ evolves according to

\begin{equation}
\tau \frac{dV(t)}{dt} = -V(t) + R\,I(t)
\label{eq:lif}
\end{equation}

\noindent where $\tau$ is the membrane time constant governing how quickly the potential leaks back toward rest, $R$ is the neuron's input resistance, and $I(t)$ is the incoming current -- in a network context, the weighted sum of incoming spikes. Whenever $V(t)$ reaches a threshold $V_{\text{th}}$, the neuron emits a spike and $V(t)$ is immediately reset to a reset potential $V_{\text{reset}}$, after which the integration in Equation~\ref{eq:lif} begins again from that reset value. The practical consequence for this project is rate coding: a stronger, more sustained input current causes $V(t)$ to reach threshold more frequently, producing a higher spike rate, and it is this spike rate -- not any single spike's individual timing -- that is taken to correspond to a ReLU unit's continuous activation value in the conversion described below.

\realfig{lif_membrane_potential_trace.png}{0.7}{A simulated LIF neuron's membrane potential over time under constant input current, showing the characteristic integrate-toward-threshold, spike, and reset pattern described by Equation~\ref{eq:lif}.}{fig:lif-trace}

The specific technique of converting an already-trained ANN into an SNN, rather than training an SNN directly, is developed in detail by \citet{rueckauer2017}, and is the theoretical basis of the \texttt{cnn2snn} toolchain used throughout this project. The conceptual mapping is that a ReLU activation's continuous output value is approximated by a LIF neuron's spike \emph{rate} over a short window of time -- a higher ReLU activation corresponds to more frequent spiking, exactly the rate-coding relationship set out in Equation~\ref{eq:lif} above.

\realfig{ann_to_snn_conversion_diagram.png}{0.8}{A conceptual before/after diagram: a ReLU unit's continuous activation value on one side, mapped to a LIF neuron's spike train over a short time window on the other, illustrating the rate-coding relationship the conversion relies on.}{fig:ann-snn-diagram}

\citet{pfeiffer2018} give an accessible bridge between the classical LIF-neuron literature and this specific deep-learning conversion use case, and note explicitly that some accuracy loss through the conversion process (typically small, single-digit percentage points) is normal and expected, a trade-off this dissertation discusses candidly in Chapter~\ref{ch:discussion} rather than treating any accuracy gap as evidence of an implementation error.

\section{Related Detection Systems}

\subsection{Shoeb (2009)}

\citet{shoeb2009} is the fair, established comparator for this project's per-patient personalisation work (Section~\ref{sec:results-c2}): a per-patient support vector machine classifier trained on each patient's own labelled seizures, reporting a mean detection latency of approximately 4.6 seconds across the patients evaluated. This is a per-patient, not cross-patient, model -- the correct comparator for a project section that is itself about the effect of a patient's own data, not for the cross-patient LOPO work in Section~\ref{sec:results-lopo}.

\subsection{Ali et al. (2024)}

\citet{ali2024} report a random-forest classifier trained on 92 hand-crafted EEG features under a genuine leave-one-patient-out evaluation protocol, achieving 72--75\% sensitivity depending on configuration. This is the fair comparator for this dissertation's LOPO cross-patient generalisation result (Section~\ref{sec:results-lopo}), since both share the true zero-shot, cross-patient evaluation design. The comparison carries an explicit caveat, stated once here so it need not be re-argued at every later reference: Ali et al.'s system uses hand-crafted features rather than a learned representation, assumes no hardware power or quantisation constraint, and -- so far as I have been able to determine -- reports no explicit equivalent of the collapse diagnostic used throughout this project to distinguish genuine detection from over-triggering artefacts. The two systems are not solving an identical problem under identical constraints, and the gap reported in Chapter~\ref{ch:results} should be read with that in mind.

\subsection{Low-Power Embedded Seizure Detection}

A separate family of prior work targets low-power embedded or wearable seizure detection directly -- systems in the spirit of BrainFuseNet, EpiDeNet, and dedicated ultra-low-power inference accelerators such as Xylo and UltraTrail. These are useful as illustrative points in the broader energy-versus-accuracy design space this dissertation occupies, but are not treated as strict head-to-head benchmarks: they typically evaluate against different datasets (frequently PEDESITE rather than CHB-MIT) under different protocols, so that a direct numerical comparison would be misleading without substantial further verification work outside this project's scope. Where referenced in Chapter~\ref{ch:discussion}, this caveat applies throughout.